% Generated draft: requires mathematical and source review.
\documentclass[11pt]{article}
\usepackage[margin=25mm]{geometry}
\usepackage{fontspec}
\setmainfont{Latin Modern Roman}
\usepackage{amsmath,amssymb,mathtools,unicode-math}
\setmathfont{Latin Modern Math}
\usepackage{enumitem,needspace,xurl,hyperref}
\setlength{\parindent}{0pt}
\setlength{\parskip}{5pt}
\setlength{\emergencystretch}{3em}
\newcommand{\N}{\mathbb{N}}
\newcommand{\Z}{\mathbb{Z}}
\newcommand{\Q}{\mathbb{Q}}
\newcommand{\R}{\mathbb{R}}
\newcommand{\C}{\mathbb{C}}
\newcommand{\F}{\mathbb{F}}
\newcommand{\A}{\mathbb{A}}
\newcommand{\PP}{\mathbb{P}}
\newcommand{\E}{\mathbb{E}}
\newcommand{\eps}{\varepsilon}
\newcommand{\dd}{\,\mathrm{d}}
\begin{document}
\section*{UnsolvedMath ID: 701306}
\textbf{Source problem code:} catalogue entry
\subsection{Definitions and mathematical statement}
Let \(G\) be a group. Use the commutator convention \([a,b]=a^{-1}b^{-1}ab\), and define iterated commutators by \([a,{}_1 b]=[a,b]\) and \([a,{}_{n+1}b]=[[a,{}_n b],b]\) for \(n\geq 1\). An element \(x\in G\) is a \emph{right \(4\)-Engel element} if
\[
\forall g\in G,\qquad [x,{}_4 g]=1.
\]
Its normal closure is
\[
N=\langle x\rangle^G=\langle g^{-1}xg\mid g\in G\rangle.
\]
For a group \(H\), define \(\gamma_1(H)=H\) and \(\gamma_{n+1}(H)=[\gamma_n(H),H]\), where the latter denotes the subgroup generated by all corresponding commutators. The group \(H\) is nilpotent if \(\gamma_{c+1}(H)=1\) for some integer \(c\geq 0\); its nilpotency class is the least such \(c\). A group is locally nilpotent if every finitely generated subgroup is nilpotent.

\begin{enumerate}
\item[(a)] For every locally nilpotent group \(G\) and every right \(4\)-Engel element \(x\in G\), must \(\langle x\rangle^G\) be nilpotent?
\item[(b)] If the answer to (a) is affirmative, is the nilpotency class of \(\langle x\rangle^G\) bounded? Under the uniform interpretation, this asks whether there exists an integer \(C\geq 0\), independent of \(G\) and \(x\), such that
\[
\gamma_{C+1}(\langle x\rangle^G)=1
\]
for every locally nilpotent group \(G\) and every right \(4\)-Engel element \(x\in G\). The supplied wording does not explicitly specify the permitted dependence of the bound; this uniform interpretation requires confirmation.
\item[(c)] For every group \(G\) and every right \(4\)-Engel element \(x\in G\), must \(\langle x\rangle^G\) be nilpotent?
\end{enumerate}
\subsection{Short English statement}
A right \(4\)-Engel element satisfies \([x,g,g,g,g]=1\) for every group element \(g\), with commutators associated to the left. The questions ask whether its normal closure is nilpotent in a locally nilpotent group, whether an affirmative answer admits a uniform class bound, and whether nilpotence holds in every group. The intended dependence of the bound in (b) is not explicit in the source.
\subsection{Sources}
\begin{itemize}
\item Local UnsolvedMath registry: permanent ID \(701306\), exact source problem code \texttt{catalogue entry}, title \emph{Kourovka Notebook Problem 17.4}. The supplied record has upstream revision \texttt{d371ee917ac9647fbb06cd6fc529deec6f3dda82} and was updated on \texttt{2026-10-16T14:00:00Z}.
\item The supplied attribution identifies A. Abdollahi as proposer, with proposal year 2010, and cites \emph{The Kourovka Notebook (Unsolved Problems in Group Theory)}, 17th issue (2010), Problem 17.4, and the 21st issue edition (2026), arXiv identifier \texttt{1401.0300v48}.
\item The supplied collection description credits E. I. Khukhro and V. D. Mazurov as editors of the Kourovka Notebook collection covering Issues 1--21 (1965--2026).
\item The registry labels this problem \emph{open}. This is supplied source metadata, not an independently verified assessment of its mathematical or literature status. No source was independently consulted for this draft.
\end{itemize}
\end{document}
